\section{Conclusions}
\label{sec:conclusions}

\begin{enumerate}[leftmargin=1.6em,itemsep=4pt,topsep=4pt,parsep=0pt]

\item No published radio technosignature search above about 30\,GHz has
reached a star within 40\,pc (\S\ref{sec:background}), and an archive
assembled for other science is at present the only way to reach one. The one
ALMA band below those searched that holds any data on these stars is
proprietary until \SbBoneRelYear{} and channelised far too coarsely to carry
a carrier search.

\item The primary experiment is \NWinA{} fine-channel windows toward
\NSysClassA{} stellar systems, covering \CvUnionA\,GHz of union bandwidth in
disconnected intervals, searched for an unresolved carrier whose frequency
drifts. A supplementary coarse-channel search for an unresolved spectral
excess adds \NWinB{} windows toward \PrSysClassB{} systems and
\CvUnionB\,GHz of union bandwidth, with no drift discrimination. Both come
from one uniform reduction of \NEB{} public execution blocks,
\TotalDataTB\,TB of raw data, and between them they span
\SurvFreqLo--\SurvFreqHi\,GHz.

\item Sensitivity is a measured completeness of one epoch and not a
threshold: $\mathrm{EIRP}_{90}$ is the power at which nine of ten carriers
injected into the real calibrated visibilities come back through the
unchanged trigger and localisation at the star, inside a single execution
block. It is the power at which a carrier would have produced a crossing,
not the power at which one would have been confirmed. The median system's
best window reaches $\EirpNinetySysMedA$\,W, with a spread across systems of
$\EirpNinetySysLoA$--$\EirpNinetySysHiA$\,W; carriers were injected into
\SensNInj{} of the \NWinA{} windows and the rest carry the class-average
factor, whose measured dispersion moves that median by \ScRsMedShiftPct{} per
cent. One bias runs one way only and is outside that calibration: an injected
carrier suffers no atmospheric decorrelation where a real one would, which
leaves the median limit optimistic by \SensDecorPct{} per cent and the
windows above \SensDecorHiGHz\,GHz by \SensDecorHiPct{} per cent.

\item All \EvNCrossRaw{} threshold crossings are reported with their
dispositions, \PrFlagXrossWord{} of them in the one execution block that
fails the data-quality criterion of Appendix~\ref{sec:blockquality} and
reported apart from the \EvNCrossQp{} in quality-passing data. Of the
\EvNCrossRaw, \MkNCoinc{} fall within $\pm\MaskHalfKms$\,km\,s$^{-1}$ of a
catalogued transition at the star's own velocity, a coincidence that is
reproducible and involves no judgement. Of those, \MkNSODecideWord{} are
with sulphur monoxide, reported and not acted on for a reason about the
species; acting on them instead would make the split \MkAttrAlt{} to
\MkUnattrAlt{} rather than \MkAttr{} to \MkUnattr. The mask costs
\MkMaskAPct{} per cent of the Class~A union bandwidth, and re-searching the
line-dominated windows with its channels removed adds no crossing. Chance,
measured from each window's own control positions and channel grid, predicts
\CeExp{} unattributed Class~A crossings over the census against the \CeObs{}
found there, or \CeExpX{} with the measured star-to-control transfer applied,
and resampling the execution blocks allows \CeLo--\CeHi. The count is
consistent with chance, and the width of that range is the limit of the
statement: an added population of as many as \CeAddBlind{} crossings would
sit inside it unremarked.

\item Every unattributed crossing was tested for recurrence in the star's
own frame, over the whole of each further window covering the event frequency
and at every drift searched, because an accelerating carrier need not return
to one channel. None was found
again: the largest statistic at any predicted cell is \RcTMax{}
against $T_{\rm trig}=\StageNullTrig$. Of the \TjNPair{} pairs of crossings
lying close enough in frequency to be each other's repeat, the two fitted
drifts and the elapsed time resolve the question for \TjNPower{} of them, and
a single carrier is excluded in every one; over the longer baselines they
resolve none, and there the criterion tests the population rather than the
pair. The same test
accepts the one set of crossings that is present again, \TjCtlStar's two
carbon monoxide crossings. Of the \RcNExcl{} crossings
that carry an exclusion, the repeat coverage rules out a carrier at a median
\RcFracMed{} of the discovery amplitude with individual limits from
\RcFracMin{} to \RcFracMax, so the median is not a bound on the set;
combining every covering block assumes a carrier that returns to one
stellar-frame channel at the start of each, and where the repeat is shallow
the criterion has little power and its silence is not an exclusion:
\RcNPowerLo{} of these crossings are untested rather than excluded. What
these tests exclude is a carrier that holds its frequency inside the window
searched and radiates through both observations, and not a persistent
transmitter of any kind.

\item What this survey could search and what it could confirm are different
extents. Of the \NSysClassA{} systems the carrier search covered,
\ScSysIndep{} were observed in at least two independently separated epochs;
for the other \ScSysNoConf, observed once or only within a single day,
confirmation was beyond reach at any sensitivity, and about those systems the
null reported here says nothing at all.

\item Combining each surveyed star's epochs in its own rest frame deepens
the flux limit by a median factor \SyStkGainMed. For Proxima~Centauri the
stacked limit, measured by injection into the stack, reaches
$\SzProxEirpCensus$\,W over the primary census alone and
$\StkProxEirpCal$\,W once the rest of the archive is added, in
\ChanBModeMHz\,MHz channels --- a Class~B limit on an unresolved excess,
carrying no drift discrimination --- for an emitter steady in the stellar
frame to within one channel across \StkProxSpanYr\,yr.

\item These results constrain unresolved carriers only within the domain
actually searched (\S\ref{sec:excludes}); what transfers beyond it is the
obstacle, that channelisation chosen by other observers, and not
collecting area, decides which kind of search is possible at all.

\end{enumerate}
